% !TeX root=../main.tex
\chapter{کالیبراسیون حسگر چند وجهی و پیاده‌سازی‌های عملی}
%\thispagestyle{empty} 
\label{chap:results}
\section{مقدمه} 
کالیبراسیون دقیق حسگرهای چندوجهی یک گام حیاتی در کاربردهای رباتیک پیشرفته، به ویژه در سناریوهایی مانند گرفتن اجسام نرم و ظریف است. با این حال، دستیابی به داده‌های کالیبراسیون با کیفیت بالا و برچسب‌زنی شده در محیط‌های واقعی با چالش‌های قابل توجهی همراه است. در این پژوهش به دلیل ماهیت پیچیده و هزینه‌بر فرآیند برچسب‌زنی، تنها مقادیر گسسته و حالت ماندگار
\LTRfootnote{Steady state}
 از نیروی عمودی و مماسی به عنوان برچسب در دسترس بود. این محدودیت، مدلسازی دینامیکی حسگر برای استخراج روابط وابسته به زمان و سناریو‌های دینامیکی را که به طور ذاتی در سیستم وجود دارند، غیرممکن می‌ساخت.

این فصل به ارائه یک رویکرد دو مرحله‌ای برای حل چالش فوق می‌پردازد. ابتدا، با استفاده از برچسب‌های استاتیک، یک مدل شبکه عصبی با معماری تماماً متصل برای تخمین اولیه نیرو آموزش داده شد. این مدل، یک سری زمانی پیوسته از مقادیر نیروی تقریبی را تولید کرد که شامل نویزهای قابل توجهی بودند. سپس، نویز خروجی این مدل اولیه به کمک فیلتر کالمن تضعیف شد.در واقع، این فیلتر به عنوان یک معلم، داده‌ها را به طور پویا صاف و نویز آن‌ها را کاهش داد تا بر اساس اطلاعات لحظات گذشته، برچسب‌های دینامیکی و با کیفیت بالا تولید شوند. در نهایت، یک شبکه عصبی حافظه‌دار با معماری حافظه طولانی کوتاه‌مدت با استفاده از داده‌های ورودی خام و برچسب‌های پاک‌سازی شده آموزش داده شد تا بتواند به طور مستقیم، نویز ورودی را حذف کرده و رفتارهای زمانی سیستم را مدل‌سازی کند.
این رویکرد ترکیبی، مزایای قابل توجهی را نسبت به روش‌های سنتی فراهم ‌آورد. در حالی که مدل استاتیک به تنهایی فقط قادر به یادگیری یک نگاشت ساده از ورودی به خروجی بود، ترکیب آن با یک شبکه عصبی حافظه‌دار، امکان پردازش داده‌های وابسته به زمان را فراهم کرد. این معماری در واقع، شروع یک امکان‌سنجی اولیه از کاربرد شبکه‌های عمیق حافظه‌دار در استخراج الگوهای پنهان در توالی داده‌ها و ترکیب منابع مختلف حسی با پهنای باند متفاوت در بستری یکپارچه بود.چنین رویکردی، توانایی ارائه‌ی خروجی دقیق‌تر در طول فرآیند تعامل با اشیاء را دارد. در این فصل، نتایج هر دو مرحله کالیبراسیون مورد بررسی قرار گرفته و عملکرد مدل‌های استاتیکی و دینامیکی با معیارهای تعریف‌شده برای دقت\LTRfootnote{Precision} و صحت \LTRfootnote{Accuracy} مقایسه می‌شود.
% در نهایت، کاربرد عملی مدل آموزش‌دیده در یک سناریوی واقعی با چنگک دو انگشتی مورد ارزیابی قرار می‌گیرد.
\section{کالیبراسیون استاتیک}
\label{sec:calib_static}
کالیبراسیون استاتیک اولین و مهم‌ترین گام برای استخراج داده‌های دقیق از حسگر است. برخلاف حسگر معرفی شده در فصل \ref{chap2}، رابطه‌ی بین نیرو و خروجی حسگر‌های اثر‌هال، یک رابطه‌ی غیرخطی بوده و این غیرخطی‌ بودن هنگام ترکیب دو حسگر اثر‌هال و یک حسگر بارومتر بیشتر می‌شود. همچنین داده‌ها برای آموزش و ارزیابی حسگر چندوجهی دارای ابعاد و پیچیدگی بیشتری نسبت به آنچه در فصل \ref{chap2} مشاهده شد، دارند. 
\subsection{مجموعه داده}
در فرآیند کالیبراسیون استاتیک، داده‌های حسگر در شرایطی جمع‌آوری شدند که مرجع برای ثبت داده‌ی واقعی در شرایط گذار نبود و تنهای در زمان‌های ماندگار حسگر، داده‌ی مرجع وجود داشته. این داده‌ها شامل چهار سری مجزا از آزمایش‌ها بودند:
\begin{enumerate}
	\item 
	اعمال نیروی عمودی ($F_z$) از صفر تا 20 نیوتون و اعمال نیروی مماسی ($F_x$) با مقادیر مشخص و تصادفی در بازه 0 تا 6 نیوتون.
	\item 
	کم کردن نیروی عمودی ($F_z$) تا صفر نیوتون و اعمال نیروی مماسی ($F_x$) در همان بازه.
	\item
	چرخاندن حسگر به اندازه‌ی 180 درجه حول محور عمود و اعمال نیروی عمودی ($F_z$) از صفر تا 20 نیوتون و اعمال نیروی مماسی ($F_x$) با مقادیر مشخص و تصادفی در بازه 0 تا $-6$ نیوتون.
	\item
	کم کردن نیروی عمودی ($F_z$) تا صفر نیوتون و اعمال نیروی مماسی ($F_x$) در همان بازه.
\end{enumerate}

   \begin{figure}[t]
	\centering
	\centerline{\includegraphics[width=0.8\textwidth]{static_nn}}
	\caption[معماری مدل استاتیک.]{
		معماری مدل استاتیک. 
	}
	\label{fig:static_architecture}
\end{figure}
 \subsection{معماری و آموزش مدل استاتیکی ارائه شده}
برای شناسایی رابطه‌ی غیرخطی،  یک مدل رگرسیون غیرخطی بر اساس شبکه‌ی عصبی تماماً متصل برای نگاشت داده‌های خام حسگر به مقادیر نیروی واقعی اعمال‌شده، توسعه داده شد. هدف این بوده که به جای استفاده از مدل‌های خطی ساده، پیچیدگی و خاصیت غیرخطی ذاتی حسگر و اثرات متقابل بین محورها در نظر گرفته شود.به عبارتی، برخلاف حسگر معرفی شده در فصل \ref{chap2} رابطه‌ی بین نیرو و خروجی حسگر‌های اثر‌هال، یک رابطه‌ی غیرخطی بوده و این غیرخطی‌ بودن هنگام ترکیب دو حسگر اثر‌هال و یک حسگر بارومتر بیشتر می‌شود. بنابر این، مدل پیشنهادی نشان داده شده در شکل
\ref{fig:static_architecture}
 ،طوری انتخاب شد که افزونگی کافی برای استخراج این رابطه ی غیر خطی را داشته باشد. 
    \begin{figure}[t]
	\centering
	\includegraphics[width=0.8\linewidth]{training_loss_static}
	\caption{خطای داده‌های آموزش و ارزیابی.}
	\label{fig:static_trainLoss}
\end{figure} 
  معماری مدل استاتیک، یک شبکه عصبی پبش خور\LTRfootnote{feed-forward} با چهار لایه‌ی تماماً متصل می‌باشد. این معماری پس از آزمون و خطا به دلیل توانایی در مدل‌سازی روابط غیرخطی پیچیده و عدم بیش‌برازش  انتخاب شده‌است. در معماری ارائه شده، لایه ورودی شامل بردار 7 تایی مقادیر حسگرهای فشار و اثر‌هال است که سوگیری اولیه‌شان
  \LTRfootnote{Bias}
   حذف شده و سپس تا بازه‌ی 0 تا یک، مقیاس‌بندی
  \LTRfootnote{Scale}
   شده‌اند. تابع فعال‌ساز همه‌ی لایه‌ها رلو
   \LTRfootnote{Relu} 
   می‌باشد. درنهایت، لایه‌ی خروجی شامل 2 نرون برای مقادیر نیروی عمودی $F_z$ و نیروی مماسی $F_x$ با تابع فعال‌ساز خطی می‌باشد که مقادیر نیرو را در بازه‌ی 0 تا 1 تخمین می‌زند؛ این مقدار با مقیاسی که بر اساس داده‌های آموزش محاسبه شده است به محدوده‌ی واقعی نگاشت می‌گردد.
  \\
  مدل بر روی مجموعه داده معرفی شده در \ref{sec:multimodal_data} آموزش داده شد. برای آموزش مدل از الگوریتم پس‌انتشار خطا
   \LTRfootnote{Error back-propagation}
   با بهینه‌ساز آدام
    \LTRfootnote{ADAM}
    برای حداقل کردن تابع هزینه‌ی میانگین مربعات خطا استفاده شد. نرخ یادگیری $0.001$ و مدل برای $200$ دور آموزش دید. نمودار خطای آموزش و ارزیابی در شکل \ref{fig:static_trainLoss} قابل مشاهده است.




\subsection{تحلیل دقت و صحت}

پس از آموزش مدل شبکه عصبی، عملکرد آن با استفاده از معیارهای دقت و صحت ارزیابی شد. معیار دقت
 میزان تکرارپذیری اندازه‌گیری‌ها را نشان می‌دهد. در یک سیستم کالیبره‌شده، اندازه‌گیری‌های تکراری تحت شرایط یکسان باید نتایج مشابهی داشته باشند. این معیار، معادل انحراف معیار مقادیر خروجی حسگر برای هر نیروی مرجع می‌باشد که بزرگتین یا همان بدترین انحراف معیار به عنوان دقت، گزارش می‌شود. در مقابل، معیار صحت میزان نزدیکی مقادیر پیش‌‌بینی‌شده به مقادیر واقعی را نشان می‌دهد. 
نمودار قدر مطلق میانگین خطا
\LTRfootnote{MAE}
و انحراف معیار خطا برای نیروی عمودی و مماسی در شکل
\ref{fig:mean_std_static}
نشان داده شده‌است. جدول 
\ref{tab:acc_prec_static}
مقادیر دقت، صحت و بازه‌‌ی کاری برای هر کدام از محور‌های نیرو را نشان می‌دهد.
    \begin{figure}[t]
	\centering
	\includegraphics[width=1.0\linewidth]{FzFx_STDandMAE_static}
	\caption{میانگین قدر مطلق خطای و انحراف معیار خطا برای  مدل استاتیک.}
	\label{fig:mean_std_static}
\end{figure} 
\begin{table}[t]
	\centering
	\caption{دقت و صحت مدل تماماً متصل.}
	\label{tab:acc_prec_static}
	\onehalfspacing
	\begin{tabular}{|r|r|r|r|}
		\hline
		\textbf{مولفه‌ی نیرو} & \textbf{صحت} & \textbf{دقت}& \textbf{بازه ‌کاری}  \\ \hline \hline
		
		$F_z$& $98.75\%$& $>99.6\%$& $0\sim20 N$ \\ \hline
		$F_x$& $98.3\%$& $>98.3\%$& $-6 \sim 6 N$ \\ \hline
		
	\end{tabular}
\end{table}
%\subsection{پسماند}
%
%هیسترزیس به پدیده وابستگی خروجی حسگر به جهت تغییر ورودی (افزایش یا کاهش نیرو) اشاره دارد. برای ارزیابی این پدیده، دو مدل شبکه عصبی مجزا آموزش داده شد: یکی برای داده‌های مربوط به اعمال نیرو (بارگذاری) و دیگری برای داده‌های مربوط به برداشتن نیرو (تخلیه بار).
%
%با استفاده از این دو مدل، یک منحنی هیسترزیس برای هر محور ترسیم شد که تفاوت بین مسیر بارگذاری و تخلیه بار را نشان می‌دهد. مساحت داخل حلقه هیسترزیس نشان‌دهنده میزان انرژی ازدست‌رفته در چرخه اندازه‌گیری و در واقع میزان پدیده هیسترزیس است. با مقایسه مقادیر پیش‌بینی‌شده توسط هر دو مدل، مشخص شد که حسگر هیسترزیس قابل‌توجهی ندارد که نشان از عملکرد پایدار آن در شرایط بارگذاری و تخلیه بار است.

\section{کالیبراسیون دینامیکی}
\subsection{رویکرد فیلتر کالمن در آماده‌سازی مجموعه داده‌ی دینامیکی}

برای غلبه بر محدودیت‌های ناشی از داده‌های استاتیک، یک رویکرد خلاقانه برای تولید برچسب‌های سری زمانی و با کیفیت بالا اتخاذ شد. در این مرحله، داده‌های حسگر خام مستقیماً به مدل استاتیک آموزش‌دیده در بخش \ref{sec:calib_static} داده شدند تا یک سری زمانی پیوسته از مقادیر نیروی تخمین‌زده شده تولید شود. این مقادیر اگرچه تخمینی از نیرو در لحظه بودند، اما به دلیل ماهیت استاتیکی مدل، نویز و نوسانات قابل توجهی داشتند.
\\
در این پژوهش، فیلتر کالمن به عنوان یک ابزار برای کاهش نویز و تولید یک مسیر صاف و عاری از نیرو مورد استفاده قرار گرفت. فیلتر کالمن یک الگوریتم بازگشتی است که به طور بهینه، حالت یک سیستم دینامیکی را بر اساس یک سری از اندازه‌گیری‌های نویزدار و یک مدل پیش‌بینی‌کننده، تخمین می‌زند. در این پژوهش، از خروجی‌های نویزدار مدل استاتیک به عنوان \textbf{اندازه‌گیری} یا\textbf{مشاهده} برای فیلتر کالمن استفاده شد. فیلتر با ترکیب هوشمندانه این مشاهدات با مدل دینامیکی ساده‌ای از حرکت سیستم (که فرض می‌کند نیروی واقعی به آرامی تغییر می‌کند)، توانست یک سری زمانی جدید از نیروهای کاهش نویز داده شده تولید کند. این سری زمانی که رفتار دینامیکی را به درستی منعکس می‌کند، به عنوان برچسب‌های دینامیکی برای آموزش مدل نهایی استفاده شدند. نحوه‌ی عملکرد فیلتر کالمن در این کاربرد به این صورت است که ابتدا یک سیستم مکانیکی ساده در نظر گرفته می‌شود:
\begin{equation}\label{eq:ss}
	\hat{x}_k = A_{k-1} \hat{x}_{k-1} + B_{k-1}u_{k-1}
\end{equation}

در این سیستم، بردار حالت $x$ مقادر کاهش نویز داده شده‌ی خروجی مدل استاتیک می‌باشند. در واقع با این فرمول بندی فرض بر آن است که تغییرات واقعی نیرو‌ها به آرامی انجام می‌شود و تغییرات سریع ناشی از نویز است. برای کاربرد کاهش نویز، ماتریس حالت $A$ به صورت همانی در نظر گرفته شده‌است. در این کاربرد فرض شده است که این سیستم ورودی ندارد پس عبارت $B_{k-1}u_{k-1}$ از \ref{eq:ss} حذف می‌شود. همچنین برای استفاده از این روش، از مقادیر خام و نویزدار خروجی مدل استاتیک به عنوان  یک بردار اندازه‌گیری یا مشاهده، $z_k$، استفاده می‌شود به همین علت ماتریس اندازه‌گیری که مشاهدات را به حالت‌ها مرتبط می‌کند، همانی خواهد بود. در ادامه می‌بایست حالات سیستم که داده‌های کاهش نویز داده شده هستند، بر اساس رابطه‌ی \ref{eq:kalman_update}، با داده‌های خام ترکیب شوند تا در عین تضعیف نویز، تغییرات را نیز دنبال کنند. 
\begin{equation}\label{eq:kalman_update}
	x_k = \hat{x}_k + K_{k}(z_k - H\hat{x}_k)
\end{equation}
در این رابطه، $H$ ماتریس اندازه‌گیری می‌باشد که طبق توضیحات، همانی در نظر گرفته شده. $K_k$ نیز بهره‌ی کالمن می‌باشد که تعیین کننده‌ی میزان اثر داده‌های اندازه‌گیری است.

اساس این روش بر مبنای وجود نایقینی در سیستم و اندازه‌گیری استوار است، نایقینی در اندازه‌گیری با ماتریس $R$ تعیین می‌شود که واریانس خطای اندازه‌گیری می‌باشد این ماتریس بیانگر قدرت نویز در خروجی خام مدل استاتیک می‌باشد که براساس مشاهدات و واریانس خروجی مدل استاتیک مقداردهی می‌شود. بر اساس مقادیر انحراف معیار گزارش شده در شکل \ref{fig:mean_std_static}  و با فرض مستقل بودن بردار‌های نیرو از هم، ماتریس $R$ به صورت زیر تعریف شد:
\begin{equation}
	R = 
	\begin{bmatrix}
		0.01&0\\
		0& 0.09
	\end{bmatrix}
\end{equation}
نایقینی در سیستم با ماتریس $P_k$ نشان داده می‌شود، این ماتریس متغیر با زمان بوده و به صورت زیر تعیین می‌گردد:
\begin{equation}\label{eq:kalman_P_update}
	P_k =AP_{k-1}A^T + Q
\end{equation}

در رابطه‌ی فوق $Q$ نماینده‌ی نایقینی‌های موجود در مدل، اعم از کم برازش یا تغییرات کوچک متغیر با زمان در مدل است. ماتریس $Q$ یک ابرپارامتر بوده و در این پژوهش به صورت زیر تعریف شده است:
\begin{equation}
	Q = 
	\begin{bmatrix}
		10^{-3}&0\\
		0& 10^{-3}
	\end{bmatrix}
\end{equation}
بر این اساس، کوواریانس تخمین در هر نمونه برداری،طبق رابطه‌ی \ref{eq:kalman_P_update} به‌روز‌رسانی می‌شود. و این مقدار در رابطه‌ی زیر، به محاسبه‌ی بهره‌ی کالمن منجر می‌گردد:
\begin{equation}\label{eq:kalman_gainupdate}
	K_k = P_kH_k\left(H_kP_kH^{T}_{k + R}\right)^{-1}
\end{equation}
در این رابطه، $R$ ماتریس کوواریانس اندازه‌گیری، $H$ ماتریس اندازه‌گیر و $P_k$ کوواریانس تخمین در آن لحظه است.
پس از اصلاح مقادیر خروجی طبق رابطه‌ی \ref{eq:kalman_update}، نایقینی تخمین طبق رابطه‌ی زیر به‌روز‌رسانی می‌شود:
 \begin{equation}\label{eq:kalman_P_update2}
 	P_k = (I - K_kH)P_k
 \end{equation}
این الگوریتم در طول سری‌های زمانی که از خروجی مدل استاتیک گرفته شد، پیاده شده و منجر به کاهش نویز این مجموعه داده گردید. برای استفاده از این داده‌ها برای آموزش و ارزیابی، از مجموعه دادگان، توالی‌های زمانی با ده نمونه‌ی متوالی استخراج شد. نحوه‌ی عملکرد مدل با این توالی‌ها به این شکل است که این ده نمونه‌ی متوالی به مدل داده می‌شود و مدل نمونه‌ی یازدهم را تخمین می‌زند. در ادامه، معماری مدل دینامیکی و نحوه‌ی آموزش و ارزیابی آن معرفی می‌شود.
\subsection{معماری و آموزش مدل دینامیکی}
   \begin{figure}[t]
	\centering
	\centerline{\includegraphics[width=0.8\textwidth]{dynamic_nn}}
	\caption[معماری مدل دینامیکی.]{
		معماری مدل دینامیکی. 
	}
	\label{fig:dynamic_architecture}
\end{figure}
    \begin{figure}[ht]
	\centering
	\includegraphics[width=0.8\linewidth]{training_loss_dynamic}
	\caption{خطای داده‌های آموزش و ارزیابی.}
	\label{fig:dynamic_trainloss}
\end{figure} 
پس از تولید برچسب‌های دینامیک، یک مدل شبکه عصبی پیچیده‌تر با قابلیت پردازش سری‌های زمانی برای کالیبراسیون نهایی طراحی شد. هدف این بود که مدل بتواند به طور مستقیم از داده‌های حسگر خام، یک نگاشت دقیق به برچسب‌های تمیز و بدون نویز پیدا کند. برای این منظور، از یک معماری ترکیبی شامل یک لایه تمام‌متصل با 24 نرون و فعال‌ساز رلو، دو لایه حافظه بلند کوتاه مدت با ابعاد پنهان
\LTRfootnote{Hidden state}
16 و 8 استفاده شد و در نهایت لایه‌ی خروجی با دو نرون، جریان داده را به نیروی عمودی و مماسی نگاشت کرد.
به علت پیچیدگی فرایند جمع‌آوری داده، از داده‌های ارتعاشی میکروفون در معماری و آموزش مدل‌ها صرف نظر شد، اما در آینده می‌توان همین مدل‌های آموزش دیده را با مدل‌های پیچیده‌تر ادغام کرد و از ماهیت چندوجهی مجموعه حسگر ساخته شده، به طور کامل بهره برد.

در نهایت، مجموعه داده‌ی جدید که متشکل از توالی‌های استخراج شده از دادگان کاهش نویز یافته بودند، برای آموزش مدل استفاده شدند. آموزش این مدل، همانند مدل استاتیک، از الگوریتم پس انتشار خطا با تابع هزینه‌ی میانگین مربعات خطا و بهینه‌ساز آدام با نرخ یادگیری $0.003$ و به تعداد 100 دور انجام شد که نمودار خطای آموزش و ارزیابی در شکل \ref{fig:dynamic_trainloss} قابل مشاهده است.



\subsection{تحلیل دقت و صحت}
همان طور که در شکل \ref{fig:dyn_static_output} نشان داده شده است، استفاده از مدل دینامیکی ارائه شده، به طور چشمگیری منجر به کاهش نوسانات خروجی گشت. برای بررسی دقت و صحت مدل جدید، از همان معیار‌های معرفی شده برای ارزیابی مدل استاتیکی استفاده شد. در این راستا، نمودار‌های میانگین و انحراف معیار خطا برای نیرو‌های عمودی و مماسی در شکل \ref{fig:mae_std_dynamic} نشان داده شده‌است. بر اساس این نمودار‌ها، دقت و صحت مدل دینامیکی در جدول
\ref{tab:acc_prec_static}
گزارش شده‌است. مدل دینامیکی تاثیر چندانی بر صحت تخمین نیرو نداشته، اما دقت تخمین را به میزان قابل توجهی افزایش داده‌است. بررسی دقیق تر داده‌ها و ارزیابی‌ها نشان می‌دهد که گلوگاه افزایش دقت و صحت مدل، داشتن داده‌های پیوسته، دقیق و کم نویز از مقدار واقعی نیرو برای کالیبره کردن مجموعه حسگر است.

\begin{figure}[t]
	\centering
	\includegraphics[width=1.0\linewidth]{FzFx_STDandMAE_dynamic}
	\caption{میانگین خطای مطلق و انحراف معیار خطا برای  مدل دینامیک.}
	\label{fig:mae_std_dynamic}
\end{figure} 
\begin{table}[t]
	\centering
	\caption{دقت و صحت مدل حافظه‌دار.}
	\label{tab:acc_prec_dyn}
	\onehalfspacing
	\begin{tabular}{|r|r|r|r|}
		\hline
		\textbf{مولفه‌ی نیرو} & \textbf{صحت} & \textbf{دقت}& \textbf{بازه ‌کاری}  \\ \hline \hline
		
		$F_z$& $98.5\%$& $>99.75\%$& $0\sim20 N$ \\ \hline
		$F_x$& $97.8\%$& $>98.1\%$& $-6 \sim 6 N$ \\ \hline
		
	\end{tabular}
\end{table}

\begin{figure}[t]
	\centering 
	\subfloat[نیروی عمودی $F_Z$.]{ 
		\includegraphics[width=0.5\textwidth]{Fz_dynamic_rise}}
	%\hspace{2mm}
	\subfloat[نیروی مماسی $F_x$.]{ 
		\includegraphics[width=0.5\textwidth]{Fx_dynamic_180rise}}%
	
	\caption{خروجی مدل استاتیکی و دینامیکی ارائه شده در مقایسه با مقادیر واقعی.}
	\label{fig:dyn_static_output}
\end{figure}
\section{پیاده‌سازی و تست‌های عملی}
پس از کالیبره کردن حسگر و شناسایی یک مدل برای نگاشت داده‌های حسگری به بردار نیرو، سیستم ساخته شده در سناریو‌های واقعی ارزیابی شد تا موثر بودن عملکرد خود را اثبات کند. در این راستا دو سناریوی کنترل نیرو و نشان دادن پدیده‌ی لغزش انتخاب شد.
\subsection{کنترل نیروی گرفتن}
\subsection{تشخیص لغزش }
به دلیل بالا بودن سطح ارتعاشات در چنگک رباتی و تحمیل اغتشاش و لرزش به مجموعه حسگر لامسه، پدیده‌ی لغزش در این حالت واضح نبود، هر چند در پژوهش‌های آتی با جمع‌آوری داده‌های دینامیکی با برچسب لغزش امکان یادگیری مدلی برای تشخیص لحظه‌ی لغزش با ترکیب داده‌های حسگر پیزوالکتریک و سایر حسگر‌ها وجود دارد. اما راه حلی جایگزین برای ارزیابی حسگر پیزوالکتریک، ترتیب داده شد. این کار در مجموعه‌ کالیبراسیون نشان داده شده در شکل \ref{fig:callib} انجام شد. ابتدا نیروی عمودی وارد بر حسگر افزایش داده شد و سپس نیروی ثابت مماسی اعمال گشت، سپس نیروی عمودی به آرامی کاهش داده شد تا صفحه‌ی بین میله‌ی اعمال نیرو و حسگر دچار لغزش شود. در تمام این مراحل، داده‌های مدل تخمین نیرو و طیف‌‌نگاره\LTRfootnote{Spectrogram} داده‌های ارتعاشی پیزوالکتریک نمایش داده شد. همان طور که در شکل \ref{fig:slip_moment} نشان داده شده است، لحظه‌ی لغزش به وضوح در نمودار طیف‌نگاره‌ی حسگر پیزوالکتریک مشهود است. 
\begin{figure}[t]
	\centering
	\includegraphics[width=0.9\linewidth]{slip_moment1}
	\caption[لحظه‌ی لغزش در نمودار طیف‌نگاره.]{
		لحظه‌ی لغزش در نمودار طیف‌نگاره.
	}
	\label{fig:slip_moment}
\end{figure} 
شکل \ref{fig:slip_moment} نشان می‌دهد که در لحظه‌ی لغزش توان برخی طیف‌های فرکانسی متعلق به حسگر پیزوالکتریک به طور مشهودی افزایش می‌یابد که به علت ارتعاش سطح تماس در لحظه‌ی لغزش می‌باشد. علاوه بر این، بررسی مقادیر نیروی عمودی و مماسی نشان می‌دهد که در آن لحظه، مقدار نیروی مماسی برای زمان بسیار کوتاهی افت می‌کند که این اتفاق به علت جدا شدن شی‌ء و حسگر در لحظه‌ی لغزش است. در این ارزیابی، از نمودار طیف‌نگاره برای نشان داده اثر لغزش بر حسگر پیزوالکتریک استفاده شد اما این بدان معنا نیست که روش های بهتر و دقیق‌تری برای این ‌کار وجود ندارد. سایر روش‌های تحلیل فرکانس مانند تبدیل موجک\LTRfootnote{Wavelet} و یا ضرایب مل
\LTRfootnote{Mel-Frequency Cepstral Coefficients (MFCC)}
می‌توانند در این تشخیص موثر باشند. همچنین، پدیده‌ی لغزش در داده‌های نیرو نیز تاثر می‌گذارد، پس می‌توان گفت که ترکیب داده‌های خام حسگرهای استاتیکی و ویژگی‌های فرکانسی استخراج شده از حسگر پیزو‌الکتریک، یکی از راه‌های تشخیص دقیق و به موقع لحظه‌ی لغزش و کنترل نیروی گرفتن است. 
%پس از کالیبراسیون، حسگر چندوجهی بر روی یک گریپر دو انگشتی نصب شد. هدف از این آزمایش، نمایش قابلیت‌های حسگر در یک سناریوی رباتیک واقعی و اندازه‌گیری نیروهای اعمال‌شده بر یک شیء در طول عملیات گرفتن بود.
%
%
%گریپر به یک شیء نرم (مانند یک اسفنج) نزدیک شده و به آن نیرو وارد می‌کند. در حین اعمال نیرو، داده‌های زیر به صورت همزمان جمع‌آوری و نمایش داده می‌شود:
%
%داده‌های حسگر نیرو/گشتاور: مقادیر نیرو در سه محور و گشتاور در سه محور که توسط مدل کالیبراسیون استاتیک و دینامیکی پردازش شده‌اند، به صورت نمودار نمایش داده می‌شود.
%
%داده‌های میکروفون ارتعاش‌سنج: ارتعاشات حسگر که ناشی از لغزش یا تغییرات ناگهانی است، توسط میکروفون ضبط شده و نمایش داده می‌شود.
%
%تحلیل نتایج:
%در طول آزمایش، مشاهده شد که با افزایش تدریجی نیروی وارد شده توسط گریپر، مقادیر نیروی حسگر نیز به صورت خطی افزایش می‌یابد. در لحظه لغزش شیء، ارتعاشات میکروفون به شدت افزایش می‌یابد که نشان از همبستگی بین داده‌های نیرو و ارتعاشات است. این همبستگی به ما اجازه می‌دهد تا یک سیستم کنترل هوشمند برای جلوگیری از لغزش ایجاد کنیم.